% Generated draft: requires mathematical and source review.
\documentclass[11pt]{article}
\usepackage[margin=25mm]{geometry}
\usepackage{fontspec}
\setmainfont{Latin Modern Roman}
\usepackage{amsmath,amssymb,mathtools,unicode-math}
\setmathfont{Latin Modern Math}
\usepackage{enumitem,needspace,xurl,hyperref}
\setlength{\parindent}{0pt}
\setlength{\parskip}{5pt}
\setlength{\emergencystretch}{3em}
\newcommand{\N}{\mathbb{N}}
\newcommand{\Z}{\mathbb{Z}}
\newcommand{\Q}{\mathbb{Q}}
\newcommand{\R}{\mathbb{R}}
\newcommand{\C}{\mathbb{C}}
\newcommand{\F}{\mathbb{F}}
\newcommand{\A}{\mathbb{A}}
\newcommand{\PP}{\mathbb{P}}
\newcommand{\E}{\mathbb{E}}
\newcommand{\eps}{\varepsilon}
\newcommand{\dd}{\,\mathrm{d}}
\begin{document}
\section*{UnsolvedMath ID: 800637}
\textbf{Source problem code:} catalogue entry
\subsection{Definitions and mathematical statement}
For a topological space \(X\), use the following standard convention: a countable Borel \(\omega\)-cover of \(X\) is a countable family \(\mathcal U\) of Borel subsets of \(X\) such that \(X\notin\mathcal U\) and every finite subset of \(X\) is contained in some member of \(\mathcal U\). Write \(\mathcal B_\Omega(X)\) for the collection of these covers. Here Borel sets are taken with respect to the topology of \(X\).

The property \(\mathsf S_{\mathrm{fin}}(\mathcal B_\Omega,\mathcal B_\Omega)\) means that for every sequence \((\mathcal U_n)_{n\in\N}\) with \(\mathcal U_n\in\mathcal B_\Omega(X)\), there exist finite subfamilies \(\mathcal F_n\subseteq\mathcal U_n\) such that
\[
\bigcup_{n\in\N}\mathcal F_n\in\mathcal B_\Omega(X).
\]
Denote this property of \(X\) by \(P(X)\).

The supplied question asks whether it is consistent that this property is closed under finite unions. More precisely, once the intended common ambient topological space \(T\) and class \(\mathscr C\) of admissible subspaces have been specified, the proposed consistency assertion is that there is a model of the intended foundational theory in which
\[
\forall k\geq 1\;\forall X_1,\ldots,X_k\in\mathscr C,\qquad
\left(\bigwedge_{i=1}^{k}P(X_i)\right)
\Longrightarrow
P\!\left(\bigcup_{i=1}^{k}X_i\right),
\]
where each space carries its subspace topology from \(T\).

The source record does not specify \(T\), \(\mathscr C\), or the foundational theory relative to which consistency is requested. It also does not supply the definitions of the cover classes. Thus the displayed formulation explains the question under the stated cover convention, but does not resolve the missing scope. In particular, choosing all sets of reals, all subspaces of a particular space, or all topological spaces would require further source evidence. The registry labels the problem open; this is a supplied classification, not an independently verified status assessment.
\subsection{Short English statement}
Can there consistently be a universe in which every finite union of spaces satisfying \(\mathsf S_{\mathrm{fin}}(\mathcal B_\Omega,\mathcal B_\Omega)\) again satisfies that property? This selection property requires choosing finitely many members from each of a sequence of countable Borel \(\omega\)-covers so that the combined chosen family is an \(\omega\)-cover. The intended class of spaces and the base theory for consistency are not specified in the supplied record.
\subsection{Sources}
\begin{itemize}
\item Supplied local UnsolvedMath registry record: permanent ID \(800637\), exact source problem code \texttt{catalogue entry}, upstream revision \texttt{d371ee917ac9647fbb06cd6fc529deec6f3dda82}. The registry attributes the problem to Tsaban, gives the year 2003, and labels it open.
\item Tsaban, \emph{Selection Principles in Mathematics: A Milestone of Open Problems} (2003), Problem 9. Supplied source URL: \url{https://arxiv.org/abs/math/0312182}. The record reports extraction by \texttt{arxiv\_latex} and an access date of 2026-10-16; these are supplied provenance details, not independent verification or a claim of access here.
\item The supplied title and background attach the reference label \texttt{[AddQuad]}. No bibliographic expansion or explanation of that label is supplied.
\end{itemize}
\end{document}
